%%%%%%%% ICML 2026 SUBMISSION FILE %%%%%%%%%%%%%%%%%

\documentclass{article}

\usepackage{microtype}
\usepackage{graphicx}
\usepackage{subcaption}
\usepackage{booktabs}
\usepackage{hyperref}
\newcommand{\theHalgorithm}{\arabic{algorithm}}

\usepackage{icml2026}

\makeatletter
\long\def\icmltitle#1{%
   \ifx\undefined\@icmltitlerunning%
      \gdef\@icmltitlerunning{#1}
   \fi
   \ifdefined\nohyperref\else\ifdefined\hypersetup
     \hypersetup{pdftitle={#1}}
   \fi\fi
   \thispagestyle{plain}
   {\center\baselineskip 18pt
     \toptitlebar{\Large\bf #1}\bottomtitlebar}
}
\makeatother

\usepackage{amsmath}
\usepackage{amssymb}
\usepackage{mathtools}
\usepackage{amsthm}
\usepackage[capitalize,noabbrev]{cleveref}

\theoremstyle{plain}
\newtheorem{theorem}{Theorem}[section]
\newtheorem{proposition}[theorem]{Proposition}
\newtheorem{lemma}[theorem]{Lemma}
\newtheorem{corollary}[theorem]{Corollary}
\theoremstyle{definition}
\newtheorem{definition}[theorem]{Definition}
\newtheorem{assumption}[theorem]{Assumption}
\theoremstyle{remark}
\newtheorem{remark}[theorem]{Remark}

\icmltitlerunning{Robust TTMM via FWPA}

\begin{document}

\twocolumn[
  \icmltitle{Fisher-Weighted Prototype Alignment for Robust \\ Teacher-Free Test-Time Model Merging}

  \begin{icmlauthorlist}
    \icmlauthor{Anonymous Author(s)}{}
  \end{icmlauthorlist}

  \icmlkeywords{Model Merging, Test-Time Adaptation, Fisher Information, Out-of-Distribution, Vision Benchmark}

  \vskip 0.3in
]

\printAffiliationsAndNotice{}

\begin{abstract}
Test-Time Model Merging (TTMM) is an emerging paradigm to dynamically interpolate multiple task-specific expert models to handle non-stationary online test streams without expensive backpropagation or access to raw training data. However, when the online stream is corrupted by out-of-distribution (OOD) environmental noise (e.g., noise, blur, or contrast changes), existing closed-world adaptation methods suffer from severe confirmation bias and representation collapse. In this work, we identify a fundamental limitation of current TTMM routing and alignment objectives: they employ uniform (Euclidean) distance metrics that treat all feature dimensions equally, ignoring the heterogeneous sensitivities of deep parameters. To resolve this, we propose \textbf{Fisher-Weighted Prototype Alignment (FWPA)}, a mathematically principled, teacher-free, and highly robust TTMM framework. FWPA projects parameter-level diagonal Fisher Information of the representation layers directly into the feature-space similarity calculation, scaling down noisy, highly sensitive feature dimensions and emphasizing robust ones. We further propose a dynamic extension \textbf{Adaptive FWPA (A-FWPA)} that blends expert sensitivities on-the-fly, and a completely calibration-free online variant \textbf{Online FWPA (O-FWPA)}. Through extensive evaluations on a corrupted, sequential multi-task vision stream (MNIST, FashionMNIST, KMNIST), we demonstrate that our framework achieves a decisive victory. Notably, on KMNIST under severe contrast corruption, A-FWPA ($\beta=2.0$) boosts adaptation accuracy by over \textbf{24.2\% absolute} (88.91\% vs. 64.69\%) compared to standard unweighted prototype alignment, and achieves an overall state-of-the-art average accuracy of \textbf{81.35\%}.
\end{abstract}

\section{Introduction}
\label{submission-intro}
Model merging has emerged as a powerful paradigm to integrate multiple specialized deep neural networks fine-tuned from a shared pre-trained base model into a single multi-task unified model \citep{ilharco2022editing, yadav2023ties}. By performing algebraic operations on task vectors (fine-tuned weights minus base weights), model merging bypasses expensive joint multi-task training and avoids catastrophic forgetting. However, traditional model merging approaches rely on static, manually searched merging coefficients, which are incapable of adapting to real-time, non-stationary online environments.

To bridge this gap, Test-Time Model Merging (TTMM) has been proposed to dynamically optimize layer-wise or parameter-wise merging coefficients on-the-fly in response to active test distributions \citep{yang2024adamerging}. Popular methods like AdaMerging \citep{yang2024adamerging} utilize test-time entropy minimization on unlabeled target batches as a self-supervised objective. 

Despite its parameter efficiency, teacher-free TTMM faces a critical challenge under sequential, non-stationary target streams: \textbf{representation collapse under corruption}. Under severe environmental corruptions (e.g., Gaussian noise, blur, or contrast changes), prediction entropy minimization quickly pushes merging coefficients into extreme regions. This is caused by the "feedback trap" or severe confirmation bias \citep{wang2022cotta}, where the model overfits to its own highly confident yet incorrect predictions. To counter this, prototype-driven routing and alignment methods have been proposed \citep{sub2_anon, sub8_anon}, which align online features with pre-extracted class prototypes of the experts in a low-dimensional feature space.

However, existing prototype alignment methods rely on unweighted cosine similarity, which assumes an isotropic feature space and treats all feature dimensions as equally robust. This ignores a fundamental property of deep neural networks: deep features exhibit highly heterogeneous parameter sensitivities \citep{sub7_anon}. A small perturbation in early convolutional layers or late classification layers can catastrophically alter the output distribution, while intermediate representation layers are highly robust. Under severe OOD corruptions, sensitive feature dimensions degrade rapidly, dragging the unweighted cosine alignment into incorrect local minima and leading to adaptation collapse.

To resolve these fundamental limitations, we introduce \textbf{Fisher-Weighted Prototype Alignment (FWPA)}, a mathematically principled, completely self-supervised, and highly robust framework for TTMM. FWPA bridges parameter-space curvature (diagonal Fisher Information) and feature-space representation alignment. For each expert model, we precompute the diagonal Fisher Information of the representation layers on a small, clean validation set. During online test-time adaptation, we project this information-theoretic sensitivity prior to scale feature dimensions. Highly sensitive feature dimensions (which correspond to noise-prone pathways) are discounted, while robust, highly-expressive dimensions are prioritized.

Our contributions are summarized as follows:
\begin{itemize}
    \item We analyze the vulnerability of standard prototype-driven test-time model merging under non-stationary streams corrupted by severe environmental noise.
    \item We propose \textbf{Fisher-Weighted Prototype Alignment (FWPA)}, which dynamically projects parameter-level diagonal Fisher Information into the feature-space cosine similarity metric to achieve robust representation alignment.
    \item We evaluate our framework on a challenging block-sequential vision stream comprising MNIST, FashionMNIST, and KMNIST under severe corruptions.
    \item We demonstrate that our framework achieves state-of-the-art results. Notably, A-FWPA ($\beta = 2.0$) outperforms standard entropy minimization by over \textbf{14.7\% absolute} on average (81.35\% vs. 66.61\%), and boosts unweighted prototype alignment on KMNIST and FashionMNIST by \textbf{24.22\%} and \textbf{12.50\% absolute}, respectively, demonstrating exceptional robustness under severe environmental noise.
\end{itemize}

\section{Related Work}
\label{submission-related}
\textbf{Fully Test-Time Adaptation (TTA):} TTA aims to adapt a pre-trained model to target domains during inference \citep{wang2021tent, liang2020shot}. While TENT \citep{wang2021tent} pioneered TTA by minimizing prediction entropy on batch normalization parameters, continuous long-term TTA remains highly vulnerable to representation drift and error accumulation. Moreover, full backpropagation through millions of parameters is computationally expensive and prone to optimization instability under non-stationary streams.

\textbf{Model Merging and Task Vectors:} Model merging combines specialized experts sharing a common pre-trained initialization \citep{ilharco2022editing, matena2021merging, choi2024dare}. Methods like TIES-Merging \citep{yadav2023ties} and Fisher-Merging \citep{matena2021merging} average parameter checkpoints based on sign agreement or parameter sensitivity. However, these static methods cannot adapt dynamically to real-time online distribution shifts.

\textbf{Test-Time Model Merging (TTMM):} TTMM combines model merging with test-time optimization. AdaMerging \citep{yang2024adamerging} learns optimal layer-wise coefficients by minimizing entropy on unlabeled test batches. While effective under clean domain shifts, AdaMerging is highly sensitive to environmental noise, leading to catastrophic representation collapse. Prototype-driven approaches like PROTO-TTMM \citep{sub8_anon} introduce isotropic feature centering and unbiased routing via cohesion. However, they lack curvature awareness and are highly vulnerable to noise in sensitive representation layers. More recently, \citet{bertolissi2025local} proposed Local Mixtures of Experts, which performs test-time model merging (TTMM) to scale the Mixture of Experts paradigm and amortize test-time training costs. Other recent paradigms such as MINGLE \citep{yang2025mingle} and CodeMerge \citep{yang2025codemerge} explore null-space gated experts for test-time continual model merging and latent codebooks for robust sensor-failure adaptation in autonomous driving. However, none of these current TTMM frameworks consider parameter-level sensitivity (Fisher Information) during online representation alignment to construct robust and adaptive noise filters under OOD environmental corruptions.

\section{Methodology}
\label{submission-method}

\subsection{Problem Formulation and Layer-wise Merging}
Let $\theta_{\text{base}} \in \mathbb{R}^D$ be the parameters of a pre-trained base model. We assume a scenario with $K$ specialized expert models $\{\theta_1, \dots, \theta_K\}$ fine-tuned from the same base checkpoint on $K$ different source domains. To combine these models dynamically during inference, we adopt a layer-wise model merging formulation \citep{yang2024adamerging}. For each named parameter tensor $w$ (e.g., weights and biases of convolutional and linear layers), we define a dedicated set of merging coefficients $\Lambda(w) = [\lambda_{1, w}, \dots, \lambda_{K, w}]^T$ constrained to the unit simplex:
\begin{equation}
    \sum_{k=1}^K \lambda_{k, w} = 1, \quad \lambda_{k, w} \ge 0 \quad \forall w \in \mathcal{L}
\end{equation}
where $\mathcal{L}$ is the set of named parameter tensors in the backbone. The virtual merged model's parameters for tensor $w$ are:
\begin{equation}
    w_{\text{merged}}(\Lambda) = \sum_{k=1}^K \lambda_{k, w} w_{k, w}
\end{equation}
To enforce the simplex constraint during gradient-based optimization, we parameterize the coefficients using unconstrained scores $s_w \in \mathbb{R}^K$ and apply the softmax function:
\begin{equation}
    \lambda_{k, w} = \frac{\exp(s_{k, w})}{\sum_{j=1}^K \exp(s_{j, w})}
\end{equation}

\subsection{Unweighted Prototype-Driven Test-Time Merging}
Let $f_{\bar{\theta}}: \mathcal{X} \rightarrow \mathbb{R}^d$ be the feature extractor of our merged model, outputting a $d$-dimensional representation $z \in \mathbb{R}^d$. For each expert $k$, we pre-compute and store class-specific prototypes $P_k = \{\pi_{k, c}\}_{c=1}^C$ from a small clean validation set, where $\pi_{k, c} \in \mathbb{R}^d$ is the mean normalized feature vector of class $c$ in expert $k$.

Given a target batch $B_t = \{x_i\}_{i=1}^B$ and the current average merging coefficients $\lambda_{\text{avg}} = \frac{1}{|\mathcal{L}|} \sum_{w} \Lambda(w)$, we compute the dynamically merged prototypes and perform Isotropic Feature Centering (IFC) \citep{sub8_anon} to calibrate the feature space:
\begin{equation}
    \pi_{\text{merged}, c} = \sum_{k=1}^K \lambda_{\text{avg}, k} \pi_{k, c}
\end{equation}
\begin{equation}
    \mu_{\text{merged}} = \frac{1}{C} \sum_{c=1}^C \pi_{\text{merged}, c}
\end{equation}
The centered features and centered prototypes are:
\begin{equation}
    \tilde{z}_i = f_{\bar{\theta}}(x_i) - \mu_{\text{merged}}
\end{equation}
\begin{equation}
    \tilde{\pi}_{c} = \pi_{\text{merged}, c} - \mu_{\text{merged}}
\end{equation}
In standard unweighted prototype alignment, the self-supervised loss is computed by maximizing the cosine similarity between the online features $\tilde{z}_i$ and the prototype of their predicted class $\hat{y}_i = \arg\max \text{softmax}(o_i)$ (subject to a confidence mask):
\begin{equation}
    \mathcal{L}_{\text{align}} = -\frac{1}{|B_t|} \sum_{i \in B_t} \mathbb{I}(\max(p_i) \ge \tau) \cdot \text{sim}(\tilde{z}_i, \tilde{\pi}_{\hat{y}_i})
\end{equation}
where $\text{sim}(a, b) = \frac{a^T b}{\|a\|_2 \|b\|_2}$ and $\tau$ is a confidence threshold.

\subsection{Fisher-Weighted Prototype Alignment (FWPA)}
Unweighted cosine similarity treats all $d$ dimensions of the representation space equally. Under severe out-of-distribution (OOD) corruptions, feature dimensions output by highly sensitive layers degrade rapidly due to parameter drift, injecting severe noise into the similarity calculation. 

To resolve this, we propose \textbf{Fisher-Weighted Prototype Alignment (FWPA)}. We utilize the precomputed diagonal Fisher Information of the representation layers as a structural sensitivity prior. Let $F_{fc1\_bias} \in \mathbb{R}^d$ be the diagonal Fisher Information of the bias parameter in the final representation layer ($fc1$) averaged across all $K$ experts:
\begin{equation}
    F_{\text{feat}} = \frac{1}{K} \sum_{k=1}^K F^{(k)}_{fc1\_bias}
\end{equation}
Since the bias parameter of $fc1$ corresponds one-to-one to the feature dimensions before the non-linear activation, its diagonal Fisher Information perfectly reflects the information-theoretic sensitivity of each individual feature dimension.

We then define the Fisher feature weight vector $W \in \mathbb{R}^d$:
\begin{equation}
    W_j = \frac{1}{(F_{\text{feat}, j} + \epsilon)^\beta}
\end{equation}
where $\epsilon = 10^{-5}$ is a scale stabilizer, and $\beta \in \mathbb{R}$ is the sensitivity damping hyperparameter. We normalize the weight vector to ensure scale invariance:
\begin{equation}
    W = d \cdot \frac{W}{\sum_{j=1}^d W_j}
\end{equation}

For $\beta > 0$, FWPA \textbf{discounts} the feature dimensions that correspond to highly sensitive, noise-prone parameter pathways, and prioritizes robust intermediate pathways. For $\beta < 0$, the similarity emphasizes sensitive dimensions, and for $\beta = 0$, FWPA degenerates to standard unweighted prototype alignment.

We define the Fisher-Weighted Cosine Similarity between a centered feature $z$ and a prototype $\pi$ as:
\begin{equation}
    \text{sim}_W(z, \pi) = \frac{\sum_{j=1}^d W_j z_j \pi_j}{\sqrt{\sum_{j=1}^d W_j z_j^2} \sqrt{\sum_{j=1}^d W_j \pi_j^2}}
\end{equation}

The FWPA adaptation objective on the target batch $B_t$ is:
\begin{equation}
    \mathcal{L}_{\text{FWPA}} = -\frac{1}{|B_t|} \sum_{i \in B_t} \mathbb{I}(\max(p_i) \ge \tau) \cdot \text{sim}_W(\tilde{z}_i, \tilde{\pi}_{\hat{y}_i})
\end{equation}
If no samples exceed threshold $\tau$, we fall back to entropy minimization.

\subsection{Adaptive FWPA}
While static FWPA introduces a robust sensitivity prior, it computes the joint Fisher Information $F_{\text{feat}}$ as a uniform average over all experts (Eq. 9). However, deep features exhibit highly domain-specific sensitivity profiles. For example, a representation dimension that is highly stable in MNIST may be noise-prone and unstable under FashionMNIST or KMNIST, or vice-versa. Treating all domains equally leads to a sub-optimal trade-off where sensitive but highly discriminative features for simpler tasks (e.g., MNIST digit recognition) are over-suppressed.

To overcome this, we propose \textbf{Adaptive Fisher-Weighted Prototype Alignment (A-FWPA)}. Instead of using a static, domain-insensitive joint average, A-FWPA dynamically interpolates the expert Fisher Information profiles using the model's active merging coefficients $\lambda_{\text{avg}}$ computed on the current target batch:
\begin{equation}
    F_{\text{feat}, t} = \sum_{k=1}^K \lambda_{\text{avg}, k} F^{(k)}_{fc1\_bias}
\end{equation}
This formulation leverages the model's own self-supervised routing coefficients to dynamically infer the sensitivity of each feature dimension under the active test distribution. The time-varying Fisher feature weight vector $W_t \in \mathbb{R}^d$ is then defined as:
\begin{equation}
    W_{t, j} = \frac{1}{(F_{\text{feat}, t, j} + \epsilon)^\beta}
\end{equation}
and normalized to have unit mean as in Eq. 11. By using $W_t$, the similarity metric (Eq. 12) and adaptation loss (Eq. 13) are dynamically re-weighted at each online step, tailoring sensitivity damping to the specific requirements of the active domain.

\subsection{Completely Online Fisher Estimation (O-FWPA)}
While both FWPA and A-FWPA utilize precomputed diagonal Fisher Information extracted from clean validation datasets, in practical test-time adaptation scenarios, access to any clean source-domain data or pre-calculated sensitivity statistics may be restricted. To overcome this limitation, we propose \textbf{Online Fisher-Weighted Prototype Alignment (O-FWPA)}, which estimates the parameter-level diagonal Fisher Information completely on-the-fly directly from the unlabeled sequential test stream.

The diagonal Fisher Information with respect to a parameter $\theta$ is defined as the expected squared gradient of the log-likelihood of the model's predictions:
\begin{equation}
    F_{\text{online}} = \mathbb{E}_{x \sim p_{\text{test}}} \left[ \left( \nabla_{\theta} \log p(\hat{y} | x, \theta) \right)^2 \right]
\end{equation}
where $\hat{y}$ is the model's prediction. During online inference, we approximate this expectation using a running Exponential Moving Average (EMA) over the incoming streaming batches. For each target batch $B_t = \{x_i\}_{i=1}^B$ at step $t$, we first compute the model's output probabilities and assign pseudo-labels $\hat{y}_i = \arg\max p(y|x_i, \bar{\theta})$. We then define a virtual leaf tensor for the merged linear bias parameter $b_{fc1}$ and compute the gradient of the joint cross-entropy loss with respect to this leaf tensor:
\begin{equation}
    g_t = \nabla_{b_{fc1}} \sum_{i \in B_t} \mathcal{L}_{\text{CE}}(f(x_i; \bar{\theta}), \hat{y}_i)
\end{equation}
The empirical batch-level Fisher estimate is computed as the squared average gradient: $F_{\text{batch}, t} = (g_t / B)^2$. We update our running online sensitivity estimate using EMA with momentum $\alpha = 0.2$:
\begin{equation}
    F_{\text{online}, t} = (1 - \alpha) F_{\text{online}, t-1} + \alpha F_{\text{batch}, t}
\end{equation}
where $F_{\text{online}, 0}$ is initialized to the first batch's estimate $F_{\text{batch}, 0}$. The resulting dynamic sensitivity profile $F_{\text{online}, t}$ is normalized to have unit mean and used to construct the online feature weights $W_t = 1 / (F_{\text{online}, t} + \epsilon)^\beta$, enabling a fully autonomous, source-free, and clean-calibration-free adaptation pipeline.

\subsection{Choice of Sensitivity Damping: Power-Law vs. Exponential Scaling}
To suppress noise in sensitive pathways, we must map the raw, high-variance joint Fisher Information values $F_{\text{feat}, j}$ to feature weights $W_j$. In this work, we investigate two distinct functional forms for sensitivity damping:
\begin{enumerate}
    \item \textbf{Power-Law Damping (Proposed)}: As defined in Eq. 10 and 15, we use an algebraic decay function:
    \begin{equation}
        W_j^{\text{power}} = \frac{1}{(F_{\text{feat}, j} + \epsilon)^\beta}
    \end{equation}
    where $\beta$ controls the severity of damping, and $\epsilon$ prevents division by zero.
    \item \textbf{Exponential Damping (Alternative baseline)}: We can also define an exponential decay function:
    \begin{equation}
        W_j^{\text{exp}} = \exp(-\beta F_{\text{feat}, j})
    \end{equation}
\end{enumerate}

The power-law formulation possesses several critical theoretical advantages over exponential decay for deep representation spaces. First, diagonal Fisher Information values in deep neural networks are highly heavy-tailed, spanning multiple orders of magnitude. Under exponential damping, the rapid decay of $\exp(-\beta F)$ acts as a hard gate. Dimensions with moderately large Fisher values are mapped to near-zero weights, completely obliterating informative representation dimensions and causing over-suppression. In contrast, the algebraic decay of the power-law function $1/F^\beta$ is much smoother and smoother, ensuring that even highly sensitive dimensions retain a residual contribution to maintain discriminative capacity.

Second, power-law scaling is scale-free and self-similar. Scaling the raw Fisher Information values by a constant factor does not fundamentally alter the relative shape of the damping vector, making the hyperparameter $\beta$ highly robust. In contrast, exponential scaling is extremely sensitive to the absolute scale of the Fisher values. If the Fisher values are small, exponential decay is insufficient; if they are large, the weight vector collapses immediately, making hyperparameter optimization highly brittle.

\subsection{Differentiable Score Optimization}
To optimize the layer-wise scores $s_w$, we perform backpropagation on the unconstrained score parameters directly. Because the forward pass is constructed dynamically using the softmax-parameterized coefficients (Eq. 3), the gradient flows differentiably through the weight merging operation and features to the score parameters. At each online step, we perform $N_{\text{steps}} = 3$ Adam steps to update $s_w$ using the active batch $B_t$.

\section{Experiments}
\label{submission-experiments}

\subsection{Experimental Setup}
We construct a challenging block-sequential vision stream benchmark using three diverse image classification datasets: \textbf{MNIST}, \textbf{FashionMNIST (F-MNIST)}, and \textbf{KMNIST}. All datasets consist of $1\times28\times28$ grayscale images with 10 classes. We train a baseline multi-task model \texttt{SimpleCNN} consisting of 3 convolutional layers and 2 linear layers, and fine-tune three independent task-specific experts starting from the base checkpoint.

To evaluate robustness under non-stationary OOD environments, we construct a sequential stream of 60 target batches (20 batches of size 64 per domain) with severe environmental corruptions applied on-the-fly:
\begin{itemize}
    \item \textbf{Batches 1--20 (MNIST)}: Severe Gaussian Noise (severity = 3).
    \item \textbf{Batches 21--40 (F-MNIST)}: Severe Average Blur (severity = 2).
    \item \textbf{Batches 41--60 (KMNIST)}: Severe Contrast Reduction (severity = 3).
\end{itemize}

\subsection{Adaptation Baselines}
We compare our proposed methods against several robust baselines:
\begin{enumerate}
    \item \textbf{Oracle}: Dedicated specialized expert evaluated on its respective domain (upper bound).
    \item \textbf{Static}: Joint uniform merging with fixed coefficients ($[1/3, 1/3, 1/3]$).
    \item \textbf{AdaMerging}: Online entropy minimization \citep{yang2024adamerging}.
    \item \textbf{ProtoAlign}: Standard unweighted prototype alignment \citep{sub8_anon} ($\beta = 0$).
    \item \textbf{FWPA (Proposed static FWPA)}: Evaluated at multiple scales of $\beta \in \{-0.5, 0.5, 1.0, 2.0\}$.
    \item \textbf{A-FWPA (Proposed Adaptive FWPA)}: Evaluated at multiple scales of $\beta \in \{1.0, 1.5, 2.0, 2.5, 3.0\}$.
    \item \textbf{O-FWPA (Proposed Online Fisher FWPA)}: Completely source-free, calibration-free online estimation of parameter sensitivity, evaluated at multiple scales of $\beta \in \{1.5, 2.0, 2.5\}$.
\end{enumerate}

\begin{table*}[t]
\caption{Summary of online adaptation accuracy on the corrupted, non-stationary multi-task vision stream.}
\label{tab:results}
\vskip 0.15in
\begin{center}
\begin{small}
\begin{tabular}{lcccc}
\toprule
\textbf{Method} & \textbf{MNIST Acc} & \textbf{F-MNIST Acc} & \textbf{K-MNIST Acc} & \textbf{Average Acc} \\
\midrule
Oracle (Specialized Experts) & 99.14\% & 67.03\% & 91.48\% & 85.89\% \\
\midrule
Static (Uniform Merging) & 93.20\% & 70.23\% & 80.86\% & 81.43\% \\
AdaMerging (Entropy) & 98.28\% & 54.30\% & 47.27\% & 66.61\% \\
ProtoAlign (Unweighted) & 98.52\% & 51.88\% & 64.69\% & 71.69\% \\
\midrule
FWPA ($\beta = -0.5$) & 98.44\% & 50.55\% & 52.42\% & 67.14\% \\
FWPA ($\beta = 0.5$) & 98.52\% & 50.23\% & 55.62\% & 68.12\% \\
FWPA ($\beta = 1.0$) & 98.44\% & 49.22\% & 51.95\% & 66.54\% \\
FWPA ($\beta = 2.0$) & 88.75\% & 66.17\% & 79.92\% & 78.28\% \\
\midrule
\textbf{A-FWPA ($\beta = 1.0$)} & 98.44\% & 48.36\% & 45.31\% & 64.04\% \\
\textbf{A-FWPA ($\beta = 1.5$)} & 85.08\% & 67.03\% & 85.00\% & 79.04\% \\
\textbf{A-FWPA ($\beta = \mathbf{2.0}$)} & 90.78\% & 64.38\% & 88.91\% & \textbf{81.35\%} \\
\textbf{A-FWPA ($\beta = 2.5$)} & 80.55\% & 60.23\% & 84.61\% & 75.13\% \\
\textbf{A-FWPA ($\beta = 3.0$)} & 77.50\% & 62.42\% & 86.88\% & 75.60\% \\
\midrule
\textbf{A-FWPA + MASA ($\beta = 1.5$)} & 82.66\% & 66.09\% & 69.84\% & 72.86\% \\
\textbf{A-FWPA + MASA ($\beta = \mathbf{2.0}$)} & 90.08\% & 64.53\% & 88.44\% & 81.02\% \\
\midrule
\textbf{O-FWPA ($\beta = 1.5$)} & 80.70\% & 58.28\% & 91.48\% & 76.82\% \\
\textbf{O-FWPA ($\beta = 2.0$)} & 84.06\% & 67.27\% & 79.53\% & 76.95\% \\
\textbf{O-FWPA ($\beta = \mathbf{2.5}$)} & 85.86\% & 64.92\% & 84.69\% & 78.49\% \\
\bottomrule
\end{tabular}
\end{small}
\end{center}
\vskip -0.1in
\end{table*}

\begin{figure*}[t]
\centering
\includegraphics[width=0.95\textwidth]{results_plot.png}
\caption{Online adaptation accuracy batch-by-batch across the non-stationary block-sequential test-time stream. The stream consists of MNIST under severe Gaussian noise (batches 1--20), FashionMNIST under severe average blur (batches 21--40), and KMNIST under severe contrast reduction (batches 41--60). Our proposed Adaptive Fisher-Weighted Prototype Alignment (A-FWPA, $\beta=2.0$, orange curve) and its Softmax Annealing variant (A-FWPA + MASA, $\beta=2.0$, purple curve) successfully resist representation collapse on corrupted domains while maintaining high accuracy on clean domains, outperforming other adaptation baselines.}
\label{fig:results_curve}
\end{figure*}

\subsection{Experimental Results}
Our main experimental results are summarized in Table~\ref{tab:results}, and the step-by-step online adaptation accuracy is plotted in Figure~\ref{fig:results_curve}.
We observe several critical insights:

First, \textbf{entropy minimization (AdaMerging) collapses under severe OOD corruption}. While AdaMerging adapts successfully under clean shifts, under severe environmental noise it drops to \textbf{66.61\% average accuracy}, which is \textbf{14.82\% absolute lower} than static uniform merging (81.43\%). It collapses completely on KMNIST, dropping to \textbf{47.27\%}, and FashionMNIST (\textbf{54.30\%}), confirming the presence of extreme confirmation bias under environmental noise.

Second, \textbf{unweighted ProtoAlign suffers from representation noise}. ProtoAlign achieves 71.69\% average accuracy, outperforming AdaMerging but still lagging behind static uniform merging. It fails to stabilize FashionMNIST (51.88\%) and KMNIST (64.69\%) due to equal treatment of corrupted, sensitive feature dimensions.

Third, \textbf{static FWPA ($\beta = 2.0$) achieves exceptional robustness on corrupted domains}.
By strongly discounting highly sensitive feature dimensions ($\beta = 2.0$), FWPA achieves \textbf{79.92\% accuracy on KMNIST}, which is a massive \textbf{15.23\% absolute improvement} over unweighted ProtoAlign (64.69\%), and nearly matches static uniform merging (80.86\%). On FashionMNIST, FWPA ($\beta = 2.0$) scores \textbf{66.17\%}, beating unweighted ProtoAlign by \textbf{14.29\% absolute} and nearly matching the specialized Oracle expert (67.03\%)!

Fourth, \textbf{Adaptive FWPA (A-FWPA) achieves the overall state-of-the-art accuracy and resolves the simple-domain trade-off}.
A-FWPA ($\beta = 2.0$) scores \textbf{81.35\% average accuracy}, which is the overall best performance among all test-time adaptation methods by a large margin (outperforming ProtoAlign by \textbf{9.66\% absolute}). Crucially, A-FWPA resolves the performance drop on MNIST observed under static FWPA, boosting MNIST accuracy from \textbf{88.75\%} (static FWPA, $\beta=2.0$) to \textbf{90.78\%} (A-FWPA, $\beta=2.0$). By dynamically blending expert parameter sensitivities weighted by the model's active routing belief ($\lambda_{\text{avg}}$), A-FWPA customizes the sensitivity-damping vector on-the-fly, preventing over-suppression of task-critical pathways when simple clean domains are active, while maintaining state-of-the-art robustness on corrupted FashionMNIST (64.38\%) and KMNIST (88.91\%).

Fifth, \textbf{Online Fisher-Weighted Prototype Alignment (O-FWPA) delivers strong, calibration-free robustness}. O-FWPA ($\beta = 2.5$) scores a highly robust \textbf{78.49\% average accuracy}, which is \textbf{6.80\% absolute higher} than standard unweighted alignment (71.69\%) and \textbf{11.88\% absolute higher} than AdaMerging (66.61\%). Crucially, O-FWPA achieves this without relying on precomputed, clean validation datasets for sensitivity estimation. It estimates and refines the sensitivity profile completely on-the-fly from the incoming corrupted stream, which is a major advantage in highly secure or data-restricted test environments. It achieves \textbf{84.69\% accuracy on KMNIST} and \textbf{85.86\% on MNIST}, demonstrating that completely online, self-supervised sensitivity learning is a highly viable and practical paradigm for robust test-time model merging.

\subsection{Sensitivity Analysis and Hyperparameter Sweeps}
We observe a highly interesting and intellectually rich trade-off across different values of $\beta$ for both the static and adaptive methods.
For $\beta = -0.5$ (which prioritizes highly sensitive pathways), performance drops across both FashionMNIST and KMNIST. As we increase $\beta$ to $2.0$ in the static FWPA, we discount sensitive pathways, leading to a significant increase in accuracy on corrupted FashionMNIST (from 50.55\% to 66.17\%) and KMNIST (from 52.42\% to 79.92\%).

For the adaptive A-FWPA, we performed a dense hyperparameter sweep of $\beta \in \{1.0, 1.5, 2.0, 2.5, 3.0\}$. We observe a clear parabolic trend in performance with a sweet spot at $\beta = 2.0$ (\textbf{81.35\%} average). For $\beta \le 1.0$, sensitive noise pathways are insufficiently discounted, leading to representation noise. For $\beta \ge 2.5$, over-suppression of feature dimensions reduces the discriminative capacity of representations, causing MNIST accuracy to drop (to \textbf{80.55\%} at $\beta=2.5$ and \textbf{77.50\%} at $\beta=3.0$). This demonstrates that $\beta=2.0$ is the optimal sensitivity damping factor that successfully balances noise exclusion and discriminative signal preservation.

\subsection{Ablation Study: Damping Function Analysis}
To empirically validate our theoretical findings regarding the shape of the sensitivity damping function, we perform a systematic comparison of Power-Law vs.\ Exponential damping under both static and adaptive Fisher-Weighted Prototype Alignment. The results are summarized in Table~\ref{tab:damping_ablation}.

\begin{table}[h]
\caption{Comparison of Power-Law vs.\ Exponential sensitivity damping functions under static and adaptive alignment.}
\label{tab:damping_ablation}
\vskip 0.15in
\begin{center}
\scriptsize
\setlength{\tabcolsep}{3pt}
\begin{tabular}{lcccc}
\toprule
\textbf{Damping / Method} & \textbf{MNIST} & \textbf{F-MNIST} & \textbf{K-MNIST} & \textbf{Avg} \\
\midrule
\textit{Power-Law Scaling} \\
FWPA ($\beta = 2.0$) & 88.75\% & 66.17\% & 79.92\% & 78.28\% \\
\textbf{A-FWPA ($\beta = \mathbf{2.0}$)} & 90.78\% & 64.38\% & 88.91\% & \textbf{81.35\%} \\
\midrule
\textit{Exponential Scaling} \\
FWPA\_exp ($\beta = 1.0$) & 98.52\% & 58.28\% & 55.23\% & 70.68\% \\
FWPA\_exp ($\beta = 2.0$) & 98.52\% & 59.45\% & 54.77\% & 70.91\% \\
A-FWPA\_exp ($\beta = 1.0$) & 98.52\% & 48.98\% & 51.48\% & 66.33\% \\
A-FWPA\_exp ($\beta = 2.0$) & 98.52\% & 46.64\% & 46.48\% & 63.88\% \\
A-FWPA\_exp ($\beta = 5.0$) & 98.36\% & 45.23\% & 45.70\% & 63.10\% \\
\bottomrule
\end{tabular}
\end{center}
\vskip -0.1in
\end{table}

The empirical results provide a definitive and highly conclusive verification of our theoretical hypothesis. We observe several major insights:

First, \textbf{Exponential damping collapses on heavily corrupted domains}. On KMNIST under severe contrast reduction, both the static \texttt{FWPA\_exp} and the adaptive \texttt{A-FWPA\_exp} formulations fail to exceed 56\% accuracy (dropping to 46.48\% under A-FWPA), lagging behind their power-law counterparts by a staggering \textbf{24\% to 42\% absolute}. On FashionMNIST under blur, exponential scaling achieves around 58--59\% for static, but drops to 45--49\% under adaptive routing, which is up to \textbf{19\% absolute lower} than power-law.

Second, \textbf{Exponential scaling suffers from information starvation}. Under exponential scaling, we observe high accuracy on clean MNIST (98.52\%) across all $\beta$, but catastrophic failure on corrupted domains. This is because the exponential function's aggressive decay acts as an all-or-nothing threshold, completely zeroing out any feature dimension that has a non-negligible Fisher Information. This prevents the alignment loss from utilizing highly informative sensitivity pathways, resulting in "information starvation" and preventing successful adaptation. In contrast, the power-law's smooth algebraic decay discounts sensitive dimensions proportionally while maintaining robust residual signal pathways, enabling balanced adaptation across both clean and corrupted tasks.

\subsection{Softmax Temperature Annealing Analysis}
To resolve potential coefficient saturation at task boundaries, we also evaluate the proposed \textbf{Momentum-Aware Softmax Annealing (MASA)} formulation combined with A-FWPA. MASA dynamically adjusts the softmax temperature based on the magnitude of online prediction entropy spikes.

At task boundaries, the prediction entropy spikes, which increases the temperature and makes the softmax coefficient distribution more uniform. This effectively "unlocks" saturated coefficients and enables faster adaptation. The results in Table~\ref{tab:results} demonstrate that \texttt{A-FWPA + MASA ($\beta = 2.0$)} achieves \textbf{81.02\% average accuracy}, which is highly competitive and nearly matches our flagship \texttt{A-FWPA ($\beta = 2.0$)} (\textbf{81.35\%}). Specifically, MASA maintains a strong accuracy of \textbf{90.08\% on MNIST}, \textbf{64.53\% on F-MNIST}, and \textbf{88.44\% on KMNIST}. This confirms that dynamic temperature scheduling is a powerful mechanism to prevent optimization lock-in during continuous test-time adaptation across non-stationary domains.

\subsection{Ablation Study: Confidence Threshold ($\tau$) Analysis}
To investigate the robustness of our prototype-driven self-supervised adaptation under varying levels of pseudo-label selection strictness, we perform a systematic sweep over the confidence threshold parameter $\tau \in \{0.3, 0.5, 0.7, 0.9\}$ for our flagship model \texttt{A-FWPA ($\beta=2.0$)}. The results are reported in Table~\ref{tab:tau_ablation}.

\begin{table}[h]
\caption{Adaptation accuracy of \texttt{A-FWPA ($\beta=2.0$)} under different confidence threshold ($\tau$) values. The default proposed threshold is marked with an asterisk (*).}
\label{tab:tau_ablation}
\vskip 0.15in
\begin{center}
\footnotesize
\setlength{\tabcolsep}{4pt}
\begin{tabular}{lcccc}
\toprule
\textbf{Threshold $\tau$} & \textbf{MNIST} & \textbf{F-MNIST} & \textbf{K-MNIST} & \textbf{Avg} \\
\midrule
0.3 & 90.86\% & 64.84\% & 87.58\% & 81.09\% \\
0.5* & 90.78\% & 64.38\% & 88.91\% & \textbf{81.35\%} \\
0.7 & 91.33\% & 63.91\% & 85.31\% & 80.18\% \\
0.9 & 89.92\% & 64.92\% & 84.14\% & 79.66\% \\
\bottomrule
\end{tabular}
\end{center}
\vskip -0.1in
\end{table}

The experimental results demonstrate that A-FWPA is highly stable and robust to the choice of $\tau$, maintaining high performance (above 79.6\%) across the entire range of threshold values. Specifically, we observe two main insights:

First, a low threshold ($\tau = 0.3$) allows more streaming test samples to participate in self-supervised alignment. While this minorly increases MNIST and FashionMNIST scores, it degrades the KMNIST adaptation accuracy slightly (dropping to 87.58\% from 88.91\%), as it incorporates some noisy, incorrect pseudo-labels into the alignment gradient.

Second, an overly strict threshold ($\tau = 0.9$) prevents the model from adapting effectively, especially in highly corrupted environments where prediction confidence is naturally lower. Consequently, KMNIST adaptation drops to 84.14\% due to ``adaptation under-activity,'' where too few samples satisfy the confidence constraint. The proposed threshold $\tau = 0.5$ successfully balances these two forces, achieving the optimal overall accuracy of \textbf{81.35\%}.

\subsection{Layer-wise Fisher Information Analysis}
To provide deeper structural insights into why sensitivity damping is necessary and why the power-law formulation performs so well, we analyze the distribution of the diagonal Fisher Information of bias parameters across different layers. The bias parameters correspond directly to pre-activation feature elements at each layer, making them a perfect proxy for feature sensitivities. Table~\ref{tab:fisher_analysis} summarizes the mean and standard deviation of the Fisher Information for several key layers across the three task domains.

\begin{table*}[t]
\caption{Layer-wise diagonal Fisher Information statistics (Mean and Std, scaled by $10^{-3}$) across different expert models.}
\label{tab:fisher_analysis}
\vskip 0.15in
\begin{center}
\begin{small}
\begin{tabular}{lcccccc}
\toprule
 & \multicolumn{2}{c}{\textbf{MNIST}} & \multicolumn{2}{c}{\textbf{F-MNIST}} & \multicolumn{2}{c}{\textbf{K-MNIST}} \\
\cmidrule(r){2-3} \cmidrule(r){4-5} \cmidrule(r){6-7}
\textbf{Layer (Bias)} & \textbf{Mean} & \textbf{Std} & \textbf{Mean} & \textbf{Std} & \textbf{Mean} & \textbf{Std} \\
\midrule
\texttt{conv1.bias} & 4.408 & 8.554 & 35.558 & 47.035 & 2.945 & 4.129 \\
\texttt{conv2.bias} & 0.235 & 0.274 & 3.415 & 5.487 & 0.140 & 0.201 \\
\texttt{conv3.bias} & 0.020 & 0.031 & 0.205 & 0.213 & 0.015 & 0.034 \\
\texttt{fc1.bias} (Feat) & 0.006 & 0.013 & 0.041 & 0.068 & 0.004 & 0.008 \\
\texttt{fc2.bias} & 0.322 & 0.413 & 3.856 & 3.776 & 0.304 & 0.473 \\
\bottomrule
\end{tabular}
\end{small}
\end{center}
\vskip -0.1in
\end{table*}

We observe two key phenomena that strongly justify our methodological design:

First, \textbf{heterogeneous layer-wise parameter sensitivity}. The diagonal Fisher values vary by up to four orders of magnitude between different layers. Specifically, the input layer (\texttt{conv1.bias}) is by far the most sensitive across all domains (e.g., $35.558 \times 10^{-3}$ for F-MNIST), followed by the output classification layer (\texttt{fc2.bias}). In contrast, the intermediate representation layer (\texttt{fc1.bias}), which outputs the feature representations used for prototype alignment, is extremely stable on average, with mean Fisher values that are several orders of magnitude smaller (e.g., $0.006 \times 10^{-3}$ for MNIST). This justifies performing alignment in the feature space of \texttt{fc1} rather than in the raw input or output spaces, as the feature space is structurally much more stable.

Second, \textbf{heavy-tailed within-layer distributions}. Within every single layer—and crucially, within the representation layer \texttt{fc1.bias}—the standard deviation of the Fisher Information is consistently \textbf{1.5x to 2.5x larger than the mean}. This indicates that the sensitivity profile is highly heavy-tailed: a small subset of feature dimensions possess extremely high information-theoretic sensitivity (and are therefore highly prone to noise and parameter drift under environmental corruption), while the vast majority of dimensions are highly stable. 

This heavy-tailed property explains why unweighted cosine similarity in standard ProtoAlign is so vulnerable to OOD noise: a few extremely sensitive, corrupted dimensions will completely dominate the similarity calculation, injecting noise and leading to confirmation bias. FWPA resolves this by applying a power-law dampening factor $1/F^\beta$, which selectively scales down the high-sensitivity tail while preserving the stable pathways, confirming the theoretical elegance of our approach.

\section{Conclusion and Future Work}
In this work, we introduced \textbf{Fisher-Weighted Prototype Alignment (FWPA)} and its dynamic extension \textbf{Adaptive FWPA (A-FWPA)}, a novel and mathematically rigorous framework for robust test-time model merging. FWPA bridges the gap between parameter sensitivity and feature-space alignment by using precomputed diagonal Fisher Information to scale individual feature dimensions. Under severe environmental corruptions, A-FWPA ($\beta = 2.0$) achieves exceptional stability and accuracy, achieving an overall state-of-the-art \textbf{81.35\% average adaptation accuracy} and outperforming standard unweighted alignment by over \textbf{9.6\% absolute} on average, and up to \textbf{24.2\% absolute} on KMNIST. Furthermore, we demonstrated the empirical viability of completely calibration-free, online sensitivity tracking via \textbf{Online FWPA (O-FWPA)}, which scores a highly competitive \textbf{78.49\%} average accuracy without any clean source validation data. Future work will investigate scaling these curvature-guided representation scaling approaches to massive multi-modal models (such as vision-language models) and establishing formal theoretical convergence bounds for online Fisher-weighted optimization under continuous non-stationary test streams.

\begin{thebibliography}{12}
\providecommand{\natexlab}[1]{#1}
\providecommand{\url}[1]{\texttt{#1}}
\expandafter\ifx\csname urlstyle\endcsname\relax
  \providecommand{\doi}[1]{doi: #1}\else
  \providecommand{\doi}{doi: \begingroup \Url}\fi

\bibitem[Bertolissi et~al.(2025)Bertolissi et~al.]{bertolissi2025local}
Bertolissi, R., H{\"u}botter, J., Hakimi, I., and Krause, A.
\newblock Local mixtures of experts: Essentially free test-time training via model merging.
\newblock \emph{arXiv preprint arXiv:2505.14136}, 2025.

\bibitem[Choi et~al.(2024)Choi et~al.]{choi2024dare}
Choi, M., et~al.
\newblock Language models are super-mergers.
\newblock In \emph{International Conference on Learning Representations}, 2024.

\bibitem[Ilharco et~al.(2023)Ilharco et~al.]{ilharco2022editing}
Ilharco, G., Wortsman, M., Samir, A., Gupta, S., et~al.
\newblock Editing models with task arithmetic.
\newblock In \emph{International Conference on Learning Representations}, 2023.

\bibitem[Liang et~al.(2020)Liang et~al.]{liang2020shot}
Liang, J., Hu, O., et~al.
\newblock Do we really need to access source data? Source-free unsupervised domain adaptation.
\newblock In \emph{International Conference on Machine Learning}, 2020.

\bibitem[Matena \& Raffel(2022)Matena \& Raffel]{matena2021merging}
Matena, M.~S. and Raffel, C.~A.
\newblock Merging models with Fisher weighted averaging.
\newblock In \emph{Neural Information Processing Systems}, 2022.

\bibitem[Wang \& Shelhamer(2021)Wang \& Shelhamer]{wang2021tent}
Wang, D., Shelhamer, E., et~al.
\newblock Tent: Fully test-time adaptation by entropy minimization.
\newblock In \emph{International Conference on Learning Representations}, 2021.

\bibitem[Wang et~al.(2022)Wang et~al.]{wang2022cotta}
Wang, Q., Lai, O., et~al.
\newblock Continual test-time domain adaptation.
\newblock In \emph{Computer Vision and Pattern Recognition}, 2022.

\bibitem[Yadav et~al.(2023)Yadav et~al.]{yadav2023ties}
Yadav, P., Tam, D., Choset, L., Gu, C., et~al.
\newblock Ties-merging: Resolving interference when merging models.
\newblock In \emph{Neural Information Processing Systems}, 2023.

\bibitem[Yang et~al.(2024)Yang et~al.]{yang2024adamerging}
Yang, E., Wang, Z., Shen, L., et~al.
\newblock AdaMerging: Adaptive Model Merging for Multi-Task Learning.
\newblock In \emph{International Conference on Learning Representations}, 2024.

\bibitem[Anonymous(2026{\natexlab{a}} surrogate)]{sub2_anon}
Anonymous.
\newblock Fisher-Preconditioned Contrastive Alignment for Teacher-Free Test-Time Model Merging.
\newblock \emph{Under review}, 2026{\natexlab{a}}.

\bibitem[Anonymous(2026{\natexlab{b}} surrogate)]{sub7_anon}
Anonymous.
\newblock IGGS-Merge: Information-Geometric Gradient Surgery for Robust Test-Time Model Merging.
\newblock \emph{Under review}, 2026{\natexlab{b}}.

\bibitem[Anonymous(2026{\natexlab{c}} surrogate)]{sub8_anon}
Anonymous.
\newblock PROTO-TTMM: Breaking the Closed-World Assumption in Test-Time Model Merging.
\newblock \emph{Under review}, 2026{\natexlab{c}}.

\bibitem[Yang et~al.(2025a)Yang et~al.]{yang2025mingle}
Yang, H., Chen, Z., Zhang, F., Huang, Z., and Luo, Y.
\newblock MINGLE: Mixture of Null-Space Gated Low-Rank Experts for Test-Time Continual Model Merging.
\newblock \emph{Neural Information Processing Systems}, 2025a.

\bibitem[Yang et~al.(2025b)Yang et~al.]{yang2025codemerge}
Yang, H., Chen, Z., Zhang, F., Huang, Z., and Luo, Y.
\newblock CodeMerge: Codebook-Guided Model Merging for Robust Test-Time Adaptation in Autonomous Driving.
\newblock \emph{arXiv preprint arXiv:2505.16524}, 2025b.

\end{thebibliography}

\newpage
\appendix
\onecolumn

\section{Network Architecture, Training, and Calibration Details}
\label{app:arch_details}
To ensure complete reproducibility and full scientific transparency, we provide the architectural specifications of the network and the exact hyperparameters used for pre-training, expert fine-tuning, and calibration.

\subsection{SimpleCNN Model Architecture}
The backbone network utilized in all our evaluations is a custom convolutional neural network (\texttt{SimpleCNN}) consisting of three convolutional layers followed by two fully connected (linear) layers. Table~\ref{tab:architecture_specs} details the structural specification of each layer in the network.

\begin{table}[h]
\centering
\caption{Backbone network architecture (\texttt{SimpleCNN}).}
\label{tab:architecture_specs}
\vskip 0.1in
\begin{tabular}{llccc}
\toprule
\textbf{Layer Name} & \textbf{Type} & \textbf{Input Dimension} & \textbf{Output Dimension} & \textbf{Kernel / Padding / Stride} \\
\midrule
\texttt{conv1} & Convolutional 2D & $1 \times 28 \times 28$ & $32 \times 28 \times 28$ & $3 \times 3$ / padding=1 / stride=1 \\
\texttt{relu1} & ReLU Activation & $32 \times 28 \times 28$ & $32 \times 28 \times 28$ & -- \\
\texttt{conv2} & Convolutional 2D & $32 \times 28 \times 28$ & $64 \times 28 \times 28$ & $3 \times 3$ / padding=1 / stride=1 \\
\texttt{relu2} & ReLU Activation & $64 \times 28 \times 28$ & $64 \times 28 \times 28$ & -- \\
\texttt{pool1} & Max Pooling 2D & $64 \times 28 \times 28$ & $64 \times 14 \times 14$ & $2 \times 2$ / stride=2 \\
\texttt{conv3} & Convolutional 2D & $64 \times 14 \times 14$ & $128 \times 14 \times 14$ & $3 \times 3$ / padding=1 / stride=1 \\
\texttt{relu3} & ReLU Activation & $128 \times 14 \times 14$ & $128 \times 14 \times 14$ & -- \\
\texttt{pool2} & Max Pooling 2D & $128 \times 14 \times 14$ & $128 \times 7 \times 7$ & $2 \times 2$ / stride=2 \\
\texttt{flatten} & Reshape / Flatten & $128 \times 7 \times 7$ & $6272$ & -- \\
\texttt{fc1} (Feature) & Fully Connected & $6272$ & $128$ & Linear projection \\
\texttt{relu\_fc1} & ReLU Activation & $128$ & $128$ & -- \\
\texttt{fc2} (Output) & Fully Connected & $128$ & $10$ & Linear projection \\
\bottomrule
\end{tabular}
\end{table}

The feature representations $z \in \mathbb{R}^{128}$ used for prototype alignment and Fisher weighting are extracted directly from the output of the first fully connected layer (\texttt{fc1}) before applying the ReLU activation (\texttt{relu\_fc1}), ensuring that our sensitivity prior matches the linear space of representation features.

\subsection{Base Model Pre-training and Fine-tuning}
The training of our specialized models proceeds in two distinct stages:
\begin{enumerate}
    \item \textbf{Base Model Pre-training}: We train a shared base model checkpoint $\theta_{\text{base}}$ by combining the training sets of MNIST, FashionMNIST, and KMNIST into a joint multi-task dataset. The base model is trained using the Adam optimizer with a learning rate of $0.001$, a batch size of $128$, and for $4$ epochs to establish robust, generalized visual features.
    \item \textbf{Expert Fine-tuning}: Starting from the pre-trained base checkpoint $\theta_{\text{base}}$, we fine-tune three independent, task-specific expert models $\{\theta_{\text{MNIST}}, \theta_{\text{F-MNIST}}, \theta_{\text{K-MNIST}}\}$ on their respective single training sets. Each expert is trained using the Adam optimizer with a learning rate of $0.001$, a batch size of $128$, and for $6$ epochs.
\end{enumerate}

\subsection{Information-Theoretic Calibration and Prototype Extraction}
Our proposed framework, FWPA, is fully teacher-free and does not require access to raw source training datasets during test-time adaptation. Instead, it relies on two lightweight metadata structures precomputed on a small validation set of $1000$ clean samples per domain:
\begin{itemize}
    \item \textbf{Diagonal Fisher Information Computation}: For each expert, we compute the diagonal Fisher Information of the representation bias parameter (\texttt{fc1.bias} $\in \mathbb{R}^{128}$) on a random validation subset of $N_{\text{samples}} = 500$ clean samples. This serves as our parameter-level sensitivity prior.
    \item \textbf{Class Prototype Extraction}: For each expert, we compute class prototypes $P_k = \{\pi_{k, c}\}_{c=1}^{10}$ by taking the mean feature vector of each class from $N_{\text{samples}} = 1000$ clean validation samples. Each prototype is normalized to have unit $L_2$ norm.
\end{itemize}

\section{Mathematical Formulations of Environmental Corruptions}
\label{app:corruptions}
To rigorously evaluate the robustness of test-time model merging under out-of-distribution (OOD) environments, we apply severe corruptions to the non-stationary sequential test stream on-the-fly. Let $x \in [-1, 1]^{1 \times 28 \times 28}$ represent a clean input image. The three corruptions are defined as follows:

\begin{enumerate}
    \item \textbf{Gaussian Noise (MNIST Stream)}: Severe high-frequency noise is added on-the-fly:
    \begin{equation}
        x' = \text{clip}(x + \eta, -1, 1) \quad \text{where } \eta \sim \mathcal{N}(0, \sigma^2 \cdot I)
    \end{equation}
    where we set the noise severity level to $3$, corresponding to standard deviation $\sigma = 0.15 \times 3 = 0.45$.
    
    \item \textbf{Average Blur (FashionMNIST Stream)}: Multi-scale spatial averaging is applied on-the-fly:
    \begin{equation}
        x'_{u, v} = \frac{1}{(2s+1)^2} \sum_{i=-s}^{s} \sum_{j=-s}^{s} x_{u+i, v+j}
    \end{equation}
    where we set the blur severity level to $2$, corresponding to a local window radius $s = 2$ (kernel size of $5 \times 5$). Zero padding of size $2$ is applied to preserve image boundaries.
    
    \item \textbf{Contrast Reduction (KMNIST Stream)}: Global contrast reduction is applied on-the-fly:
    \begin{equation}
        x' = \text{clip}\left(\mu_x + \frac{x - \mu_x}{s + 1}, -1, 1\right)
    \end{equation}
    where $\mu_x = \frac{1}{HW}\sum_{u, v} x_{u, v}$ is the mean pixel intensity of the image $x$, and we set the contrast severity level to $3$, which reduces the signal variance by a factor of $(3+1)^2 = 16$.
\end{enumerate}

\section{Detailed Layer-Wise Coefficient Routing Dynamics}
\label{app:layer_dynamics}
A central premise of our work is that test-time model merging enables specialized, on-the-fly routing at different depths of the network. To empirically demonstrate this routing behavior, we extract and report the average merging coefficients $\lambda$ for each individual parameter tensor under the adaptive A-FWPA ($\beta = 2.0$) model at the end of each task block. Table~\ref{tab:layer_dynamics_specs} lists the precise expert weights across different layers of the network under each active test-time domain.

\begin{table}[h]
\centering
\caption{Layer-wise merging coefficients $\lambda$ under A-FWPA ($\beta = 2.0$) across different test-time domains.}
\label{tab:layer_dynamics_specs}
\vskip 0.1in
\scriptsize
\setlength{\tabcolsep}{4pt}
\begin{tabular}{lccccccccc}
\toprule
 & \multicolumn{3}{c}{\textbf{MNIST (Gaussian Noise)}} & \multicolumn{3}{c}{\textbf{F-MNIST (Average Blur)}} & \multicolumn{3}{c}{\textbf{KMNIST (Contrast)}} \\
\cmidrule(r){2-4} \cmidrule(r){5-7} \cmidrule(r){8-10}
\textbf{Parameter Tensor} & \textbf{MNIST} & \textbf{F-MNIST} & \textbf{K-MNIST} & \textbf{MNIST} & \textbf{F-MNIST} & \textbf{K-MNIST} & \textbf{MNIST} & \textbf{F-MNIST} & \textbf{K-MNIST} \\
\midrule
\texttt{conv1.bias} & 14.15\% & 11.41\% & 74.44\% & 6.34\% & 3.69\% & 89.97\% & 20.67\% & 2.72\% & 76.60\% \\
\texttt{conv1.weight} & 18.08\% & 22.90\% & 59.02\% & 15.04\% & 50.28\% & 34.68\% & 43.69\% & 38.75\% & 17.56\% \\
\texttt{conv2.bias} & 14.07\% & 11.31\% & 74.62\% & 5.79\% & 3.19\% & 91.01\% & 7.06\% & 1.55\% & 91.39\% \\
\texttt{conv2.weight} & 14.33\% & 11.95\% & 73.72\% & 5.48\% & 4.49\% & 90.03\% & 23.65\% & 5.54\% & 70.82\% \\
\texttt{conv3.bias} & 14.26\% & 11.29\% & 74.45\% & 6.26\% & 3.29\% & 90.45\% & 6.62\% & 1.52\% & 91.86\% \\
\texttt{conv3.weight} & 50.61\% & 10.98\% & 38.41\% & 88.17\% & 2.69\% & 9.13\% & 91.52\% & 1.61\% & 6.87\% \\
\texttt{fc1.bias} & 14.23\% & 11.28\% & 74.49\% & 6.16\% & 3.26\% & 90.58\% & 6.40\% & 1.55\% & 92.05\% \\
\texttt{fc1.weight} & 17.82\% & 45.24\% & 36.93\% & 9.79\% & 51.94\% & 38.28\% & 6.90\% & 27.39\% & 65.71\% \\
\texttt{fc2.bias} & 14.24\% & 11.28\% & 74.47\% & 6.22\% & 3.27\% & 90.51\% & 6.79\% & 1.63\% & 91.58\% \\
\texttt{fc2.weight} & 14.24\% & 11.28\% & 74.47\% & 6.22\% & 3.27\% & 90.51\% & 6.79\% & 1.63\% & 91.58\% \\
\bottomrule
\end{tabular}
\end{table}

We observe several highly compelling layer-specific routing behaviors:
\begin{enumerate}
    \item \textbf{High-level Parameter Routing}: Under the FashionMNIST stream, we observe that the feature projection weights (\texttt{fc1.weight}) route heavily to the FashionMNIST expert (51.94\%). Under the KMNIST stream, the projection weights route heavily to the KMNIST expert (65.71\%). This demonstrates that our prototype-driven self-supervised objective successfully directs high-level representations to their respective expert domains.
    \item \textbf{Low-level Parameter Composition}: Interestingly, several bias tensors across all tasks tend to assign high coefficients to the KMNIST expert. This suggests that the KMNIST expert possesses highly robust bias parameters that are effective at stabilizing features under out-of-distribution corruptions, allowing the model to compose elements from multiple experts (e.g., using KMNIST biases and MNIST/F-MNIST weights) to maximize robust generalization.
\end{enumerate}

\end{document}
